% --- Section: Related Work ---
% Drafted Phase 6 — values pulled from cleanup/04_phase5_data_pulled.md
% Target: ~400 words.

\section{Literature Review}
\label{sec:related}

\textbf{LLM benchmarks for quantitative reasoning.}
StatEval~\cite{lu2025stateval} catalogues 13{,}817 problems across
descriptive statistics, probability, inference, and regression,
scored mainly by answer correctness.
MathEval~\cite{liu2025matheval} compares 23 LLMs across 19
mathematical-reasoning datasets, again on accuracy. Nagarkar et
al.~\cite{nagarkar2026canllm} compare LLMs to human respondents on
statistical tasks. All three report how often the model lands on
the right number, but none asks whether the reasoning that produced
the number is well-formed or stable under perturbation.

\textbf{Bayesian-reasoning evaluation.} Longjohn et
al.~\cite{longjohn2025bayesian} propose a Bayesian framework for
aggregating benchmark results probabilistically, not for evaluating
Bayesian reasoning itself. Du et al.~\cite{du2025icecream} document
assumption-blindness in causal reasoning: LLMs produce
``convincing yet misleading'' answers without articulating the
assumptions those answers depend on. The 44.3\% missing-assumption
rate we report extends the same pattern to inferential statistics.

\textbf{Calibration.} Guo et al.~\cite{guo2017calibration}
introduced Expected Calibration Error (ECE) and showed that modern
neural networks are systematically miscalibrated despite high
accuracy. ECE is well-studied for classifiers; its application to
free-form LLM reasoning is recent and less settled. FermiEval and
related work~\cite{fermieval2025,multianswer2026,nagarkar2026canllm}
adapt calibration to natural-language outputs, but typically on a
single dimension. We report binary-pass ECE and a per-NMACR-dimension
ECE.

\textbf{LLM-as-judge.} Zheng et al.~\cite{zheng2023judge}
popularised LLM-as-judge with MT-Bench, showing that strong LLMs
can substitute for human raters on open-ended tasks. Follow-up work
documents design-choice sensitivity and self-preference
bias~\cite{park2025judge,judgment2025noise}. Brittlebench and
ReasonBench~\cite{brittlebench2026,reasonbench2025} document prompt
sensitivity and reasoning instability under perturbation. Our
protocol constrains the judge to a single dimension
(assumption-compliance) and uses an out-of-family open-weights
model (Llama 3.3 70B Turbo). Both choices address failure modes
that Feuer et al.~\cite{judgment2025noise} document.

Across these threads, evaluation usually collapses to one accuracy
or pass-rate number. We pull the three threads together into a
tri-axis ranking and report the disagreement across the same five
LLMs on the same task suite.

\textbf{Positioning.} Prior work treats accuracy, robustness, and
calibration as separable evaluation concerns and tackles them in
isolation. Quantitative-reasoning benchmarks measure accuracy at
scale but rarely report robustness or
calibration~\cite{lu2025stateval,liu2025matheval,nagarkar2026canllm}.
Calibration studies operate on classification or short-answer
formats, not on the free-form numerical and prose responses typical
of Bayesian reasoning~\cite{guo2017calibration}. LLM-as-judge work
has matured into a default validation tool, but it is usually
applied as a holistic preference signal that carries documented
biases~\cite{zheng2023judge,park2025judge,judgment2025noise}. We
are not aware of a benchmark for Bayesian or inferential
statistical reasoning that combines (i) a perturbation taxonomy
for robustness, (ii) an external single-dimension judge for
assumption-compliance, and (iii) calibration measured on the same
task set. Our contribution is that combination; it surfaces ranking
divergences a single-metric leaderboard cannot.
